\section{Related Work}
\label{sec:related}

\subsection{Multimodal hallucination evaluation and benchmark design}

Evaluating hallucination in vision--language models has crystallized around complementary protocols. \textbf{CHAIR}~\cite{rohrbach2018chair} scores whether objects mentioned in generated captions are licensed by the image, highlighting that fluent text can still be visually unsupported. \textbf{POPE}~\cite{li2023pope} recasts assessment as balanced yes/no existence probes, making object hallucination measurable under controlled prompts and revealing sensitivity to language priors. \textbf{AMBER}~\cite{amber2023} broadens the lens to multiple error types, including attributes and relations, so that ``object-only'' scores are not mistaken for full multimodal faithfulness. These instruments answer different questions: CHAIR suits open-ended captioning; POPE suits discriminative probing; AMBER stresses multi-faceted failure modes. They should not be treated as interchangeable drop-in replacements for one another.

Large-scale discriminative suites extend the same agenda with richer task and context structure. \textbf{PhD}~\cite{liu2025phd} provides a sizable VQA-style benchmark with multiple recognition tasks and adversarial prompting modes, stressing robustness of MLLMs to visually misleading or prior-triggering questions. We use CHAIR and POPE as primary reporting axes (Section~\ref{sec:exp}) while acknowledging such broader evaluation lines as context for future stress-testing.

\subsection{Decode-time mitigation and control}

A prominent class of methods intervenes \textbf{during} autoregressive decoding without updating weights. \textbf{Visual contrastive decoding (VCD)}~\cite{leng2024vcd} contrasts logits under intact versus perturbed visual inputs to down-weight language-prior-driven objects. \textbf{OPERA}~\cite{huang2024opera} shapes decoding using attention- and retrospection-based penalties over generation trajectories. \textbf{DoLa}~\cite{chuang2024dola}, developed for text-only LMs, improves factuality by contrasting early and late layer logits; adapted to MLLMs, it remains a strong \textbf{internal} layerwise baseline but does not, by itself, enforce object-level grounding in an external visual space.

Recent MLLM-specific decoding refinements still predominantly \textbf{read or reshape internal representations}. For example, \textbf{DeCo}~\cite{wang2024deco} dynamically blends logits from earlier layers when intermediate activations appear to retain visual evidence that final layers suppress. Such designs highlight that useful visual signals often exist inside the stack, yet they inherit the universality limitations of any approach that requires access to host layer outputs and geometry.

\textbf{CHORD} differs along two axes emphasized in Section~\ref{sec:method}: (i)~it never consumes host hidden states or attention tensors---only host logits on a short candidate frontier plus deterministic tokenizer lookahead; (ii)~it scores \textbf{candidate spans} against a \textbf{frozen} region-level visual bank and a \textbf{frozen} public text encoder, so the intervention signal is defined in a shared external space. This targets the cross-modal \textbf{granularity} gap between fragmented sub-word proposals and object-level evidence, while continuous Top-$1$ monitoring limits the cost of span-level scoring.

\subsection{Post-hoc correction, preference optimization, and position relative to CHORD}

Complementary families edit outputs \textbf{after} generation or change the model itself. \textbf{LURE}~\cite{zhou2024lure} performs lightweight post-hoc revision informed by statistical factors linked to object hallucination. \textbf{Woodpecker}~\cite{yin2024woodpecker} runs a multi-stage, tool-assisted pipeline that extracts claims, queries external vision modules, and rewrites inconsistent text. \textbf{Volcano}~\cite{lee2024volcano} iterates self-feedback and revision. \textbf{HA-DPO}~\cite{zhao2023hadpo} reduces hallucination via preference optimization, permanently shifting model behavior. These methods can exploit richer verification or stronger supervision precisely because they are not constrained to the per-step budget of greedy or beam decoding.

\textbf{CHORD} targets a narrower regime: \textbf{training-free}, \textbf{in-decoding} calibration of the top-$K$ frontier using external object-level evidence and a closed-form resonance delta on lifted spans. It is closest in spirit to decode-time baselines~\cite{leng2024vcd,huang2024opera,chuang2024dola,wang2024deco} on the time axis, but closest in \emph{evidence} to tool-augmented work~\cite{yin2024woodpecker} while avoiding full post-hoc pipelines and without preference updates~\cite{zhao2023hadpo}. Section~\ref{sec:method} details how this design relates to the universality boundary and to mitigating overconfident hallucinations without relying solely on entropy triggers.
